\section{The outer questions and the permutation test}
\label{sec:game}

We now use the latent gadget to construct a permutation instance.  Its
vertices encode affine tables on tuples of parity equations.  We identify
two tables exactly when they induce the same affine map on the same retained
coordinates, including the same intercept.  This common representation will
allow us to compare an equation tuple with a sparsely projected tuple in
Section~\ref{sec:matrix}.

\subsection{Weighted parity equations and distinct positions}

A \emph{weighted parity instance} consists of a finite variable set
\(\mathcal X\), a finite nonempty set \(\mathcal E\) of equation
occurrences, and positive rational weights \(\omega_e\) summing to one.
An occurrence is an ordered equation
\[
 x_{v_{e,1}}+x_{v_{e,2}}+x_{v_{e,3}}=b_e,
 \qquad b_e\in\F_2.
\]
Its value is the weighted fraction of equations satisfied, and
\(\OPT(\mathcal E)\) denotes the maximum value.  Different occurrences
have different IDs even if their ordered equations coincide.

\begin{theorem}[Weighted parity gap]
\label{thm:parity}
For every fixed rational \(\xi>0\), there is a deterministic
polynomial-time reduction from Boolean satisfiability to weighted parity
instances, with binary-encoded rational weights, such that
\[
 \begin{array}{ll}
 \textup{YES:}&\OPT(\mathcal E)\ge 1-\xi,\\
 \textup{NO:}&\OPT(\mathcal E)\le \frac12+\xi.
 \end{array}
\]
The running-time polynomial may depend on \(\xi\).
\end{theorem}

This is the gap consequence of H\aa stad's construction in
\cite[Theorem~5.4 and its proof]{Hastad01}.  More specifically, the proof
uses a fixed dyadic verifier error \(\delta>0\), gives completeness
\(1-\delta\) and soundness \((1+\delta)/2\), and enumerates verifier
choices as equations with their rational probabilities as weights.
Choosing \(\delta\le\xi\) gives the displayed statement.  The
enumeration is a deterministic reduction; the verifier's random choices
describe the resulting weighted instance.

\begin{lemma}[Making positions distinct]
\label{lem:distinct-parity}
From a weighted parity instance \(\mathcal E\) one can deterministically
construct, in polynomial time, an instance \(\mathcal E'\) in which the
three variable IDs of every occurrence are distinct.  Its weights are
positive rationals of polynomial bit length, and
\[
 \OPT(\mathcal E')\ge\OPT(\mathcal E),
 \qquad
 \OPT(\mathcal E')
 \le \frac{1+\OPT(\mathcal E)}2+\frac3{100}.
\]
In particular, Theorem~\ref{thm:parity} with \(\xi<1/100\), followed by
this preprocessing, gives completeness at least \(1-\xi\) and NO
optimum at most \(4/5\).
\end{lemma}

\begin{proof}
Replace each variable \(v\) by one hundred clones \((v,c)\),
\(c\in[100]\).  For each original occurrence, choose its three clones
uniformly subject to choosing distinct clones whenever positions have
the same underlying variable.  List all the resulting occurrences and
give them the corresponding conditional weights.  There are at most
\(100^3\) choices per original occurrence.  Giving all clones their
original variable's label proves the completeness assertion.

For the upper bound fix arbitrary clone labels, and let
\[
 t_v=\frac1{100}\sum_{c=1}^{100}(-1)^{x_{(v,c)}}.
\]
First choose clones independently, without conditioning on distinctness.
The satisfaction probability of occurrence \(e\) is
\[
 \frac{1+(-1)^{b_e}t_{v_{e,1}}t_{v_{e,2}}t_{v_{e,3}}}{2}.
\]
Round each bias to a majority label, resolving zero biases arbitrarily.
For every equation violated by this rounded assignment, the signed
product in the numerator is nonpositive, so its satisfaction probability
is at most \(1/2\).  The other equations have probability at most one.
The weighted mean is therefore at most
\((1+\OPT(\mathcal E))/2\).  The probability of a forbidden collision
among the independently chosen clones is at most \(3/100\), by a union
bound over the three pairs of positions.  Conditioning on its complement
changes the expectation of any \([0,1]\)-valued function by at most
\(3/100\).  This proves the bound.  All conditional probabilities have
constant-size denominators apart from their original occurrence weights.
\end{proof}

For the rest of the construction, \(\mathcal E\) denotes the
preprocessed instance, so every equation has three distinct variable
IDs.  Its \emph{ordinary equation-versus-variable game} samples
\(e\) according to \(\omega\) and \(j\) uniformly from \([3]\).
The first prover receives the occurrence ID \(e\) and answers with a
satisfying triple for that equation; the second receives the variable ID
\(v_{e,j}\) and answers with a bit.  Acceptance means equality at
position \(j\).  Because the IDs within an equation are distinct,
the question pair determines the projection constraint.

For any second-prover strategy, its bits constitute a global assignment.
On an equation violated by that assignment, every valid first-prover
answer disagrees in at least one position.  Hence this ordinary
projection game, denoted \(G_{\mathcal E}\), satisfies
\begin{equation}
 \val(G_{\mathcal E})
 \le 1-\frac{1-\OPT(\mathcal E)}3
 \le \frac{14}{15}
 \quad\text{if }\OPT(\mathcal E)\le\frac45.
 \label{eq:base-game-gap}
\end{equation}
The answer alphabets have sizes four and two, respectively.

\subsection{Tuples, local answer copies, and homogeneous coordinates}

Fix a positive integer \(k\).  An equation question
\(U=(e_1,\ldots,e_k)\) is sampled from \(\omega^{\otimes k}\), with
replacement.  Its answer space is the product of the \(k\) affine
solution planes.  Blocks always use local copies: even if the same
equation or variable ID appears in several blocks, their local answers
are separate coordinates.

For a parameter \(0<\beta<1\), independently in every block either
retain the equation, with probability \(1-\beta\), or retain just the
variable ID at a uniformly chosen position, with probability \(\beta\).
The resulting question is denoted \(O\), and \(\pi\) is the
coordinate projection from answers on \(U\) to answers on \(O\).
The question \(O\) records an equation ID or a variable ID at every
ordered block position.  A singleton block accepts either bit.

Represent an equation answer by its first two bits \((x_i,y_i)\);
the third is \(b_{e_i}+x_i+y_i\).  A singleton is represented by its
actual bit.  Add a common homogeneous coordinate \(t_0\).  Thus
\[
 L_U=\F_2^{1+2k},\qquad m=1+2k,
\]
and \(L_O\) has one homogeneous coordinate, two coordinates for each
retained equation, and one for each singleton.  In either space, vectors
whose first coordinate is one are exactly valid answers.  Write
\(e_0^\top\) for the functional reading that coordinate.

The affine projection \(\pi\) extends to a linear surjection
\(\bar\pi:L_U\to L_O\) preserving \(t_0\).  On retained equation
blocks it is the identity, and on a block projected to a singleton its
coordinate is
\begin{equation}
 \begin{cases}
 x_i,&\text{position }1,\\
 y_i,&\text{position }2,\\
 b_{e_i}t_0+x_i+y_i,&\text{position }3.
 \end{cases}
 \label{eq:outer-projection}
\end{equation}
We call every map obtained by these blockwise choices a \emph{permitted
projection}; this class of maps is independent of \(\beta\).
Surjectivity follows by fixing \(t_0\) and choosing coordinates
independently in each block.  We use the same conventions for questions
with some blocks omitted.

\subsection{Exact table keys and folding}

The identification of tables across projected questions follows the
virtual-table approach of Khot--Safra~\cite[Section~1.1]{KS13}.
We first express each table on the coordinates through which it factors.
This will give the same representation to a projected table and its
pullback.

Fix latent data \((\mathcal V,K,C,\mu)\) from
Lemma~\ref{lem:latent}, where \(K=\F_2^\ell\), and fix a linear
complement \(\mathcal V=K\oplus W\).  A \emph{table} on a question
\(Q\) is an affine map from its answer space to \(\mathcal V\), or
equivalently its unique linear extension
\(P\in\Hom(L_Q,\mathcal V)\).  Here \(Q\) may be an equation
question or one of its permitted projections, and the same definition
applies to a question with omitted blocks.  Write \(c\) for the
coefficient of the homogeneous coordinate \(t_0\) in \(P\).

In an equation block, inspect the rank of its slope map
\(\F_2^2\to\mathcal V\), whose contribution we write as
\(a_i x_i+d_i y_i\).  If its rank is two, retain the equation
occurrence and this contribution.  If its rank is zero, omit the block.
In the remaining case the rank is one, so the coefficients have exactly
one of the forms
\[
 (a_i,d_i)=(v,0),\quad (0,v),\quad (v,v),
 \qquad v\ne0.
\]
These forms retain only the variable at position one, two, or three,
respectively, with contribution \(v z_i\) in its actual bit \(z_i\).
For the third form,
\[
 v(x_i+y_i)=v z_i+b_{e_i}v t_0,
 \qquad z_i=b_{e_i}t_0+x_i+y_i,
\]
so this replacement also adds \(b_{e_i}v\) to the homogeneous
intercept.  In a singleton block, retain its variable ID and its original
contribution exactly when its slope is nonzero; otherwise omit the block.

The retained blocks form the \emph{support} of \((Q,P)\).  They keep
their indices in \([k]\), their equation/variable tags, and their IDs.
A retained equation keeps its ordered first two coordinates, and a retained
singleton uses its actual-bit coordinate.  Thus repeated IDs in different
blocks still have separate local coordinates.  The support may be empty;
its homogeneous coordinate remains in every case.

Let \(J_3\) be the set of equation blocks reduced to their third bit,
and write \(v_i\) for the nonzero slope in such a block.  The
\emph{reduced table} \(\widetilde P\) on the support has the retained
block contributions specified above and homogeneous intercept
\begin{equation}
 \widetilde c=c+\sum_{i\in J_3} b_{e_i}v_i.
 \label{eq:reduced-table-intercept}
\end{equation}
To check the complete table, let \(L_{\rm supp}\) be the homogeneous
answer space of the support and let
\(\rho_{Q,P}:L_Q\to L_{\rm supp}\) retain \(t_0\) and the
specified block coordinates.  It is the identity on retained equations
and on retained singleton blocks already present in \(Q\), uses the
corresponding actual-bit coordinate from
Equation~\eqref{eq:outer-projection} on a reduced equation, and discards
omitted blocks.  Choosing the coordinates independently in each
block shows that \(\rho_{Q,P}\) is surjective.  The block identities
and Equation~\eqref{eq:reduced-table-intercept} give
\[
 P=\widetilde P\rho_{Q,P}.
\]
Surjectivity makes \(\widetilde P\) the unique linear map with this
factorization.  We define the \emph{key} of \((Q,P)\) to be its
support together with this entire reduced table.  Thus two tables have
the same key exactly when their retained coordinates agree, with the same
positions, tags, and IDs, and their induced affine maps agree on those
coordinates, including their intercepts.

We next check compatibility with every permitted projection
\(\bar\pi:L_U\to L_O\).  Let
\(P'\in\Hom(L_O,\mathcal V)\).  An equation retained in \(O\)
has the same contribution in \(P'\) and its pullback \(P'\bar\pi\),
so its reduction is identical.  A singleton of slope zero disappears on
both sides.  A singleton of nonzero slope \(v\) pulls back to the
corresponding one of the three rank-one patterns above, which reduces to
the same variable ID at the same block position.  For a third-position
singleton, the pullback adds \(b_{e_i}v\) to the intercept, and its
reduction adds \(b_{e_i}v\) once more.  The two terms cancel over
\(\F_2\).  All other intercept corrections come from retained equation
blocks and are identical on the two sides.  The supports and complete
reduced tables therefore agree: \((U,P'\bar\pi)\) and \((O,P')\)
have the same key.

Define the key universe to consist of the keys of all tables on equation
questions \(U\in\mathcal E^k\).  The compatibility just proved shows
that the key of every table on a permitted projected question already
belongs to this universe, by pullback from any compatible \(U\).

We now apply the folding principle of Bellare, Goldreich, and
Sudan~\cite[Section~3.3]{BGS98}: values on one representative determine
the values on its translates.  Adding a constant \(h\in K\) to a key
means adding \(h\) to its reduced intercept, leaving its support and
slopes fixed.  This is realized by replacing a presenting table \(P\)
with \(P+h e_0^\top\), so it preserves the key universe.  The action
is free because the key records the intercept.  A vertex of our
permutation instance is an orbit of this action, represented by its unique
key whose reduced intercept has zero \(K\) component.

For a table \(P\) on an equation question or a permitted projected
question \(Q\), let \(c(Q,P)\in K\) be the \(K\) component of
its key's reduced intercept.  Let \(v(Q,P)\) denote the orbit vertex
represented by the key obtained after removing this component.
A labeling \(\lambda\) of the vertices defines the restored table answer
\[
 F_Q(P)=\lambda(v(Q,P))+c(Q,P).
\]
Adding a constant changes only the removed component, so
\begin{equation}
 F_Q(P+h e_0^\top)=F_Q(P)+h
 \qquad(h\in K).
 \label{eq:table-folding}
\end{equation}
Tables with the same key have the same representative and removed
component.  The pullback compatibility therefore gives, for every
permitted \(U,O,\bar\pi\) and every table \(P'\) on \(O\),
\begin{equation}
 F_U(P'\bar\pi)=F_O(P').
 \label{eq:key-sharing}
\end{equation}
Projected questions thus use answers determined by the same vertex
labeling, without an additional consistency test.

\subsection{The latent matrix test and completeness}

The permutation instance, denoted \(\mathcal G\), samples independently
\[
 U\sim\omega^{\otimes k},\qquad
 P\text{ uniform in }\Hom(L_U,\mathcal V),\qquad
 a\sim\mu,\qquad l\text{ uniform in }\F_2^m,
\]
and checks
\begin{equation}
 F_U(P)=F_U(P+a l^\top).
 \label{eq:matrix-test}
\end{equation}
Here \(a l^\top\) denotes the map \(z\mapsto a(l^\top z)\).
Each sample specifies an edge between the two key-orbit vertices.  If
their restored offsets are \(c_0,c_1\), its constraint is the
permutation
\[
 \lambda(v_1)=\lambda(v_0)+c_0+c_1
\]
of the fixed alphabet \(K\).  Edge occurrences have their sampling
probabilities as weights.  At this stage loops and repeated edges are
allowed.  The sparse projection experiment involving \(O\) is used to
analyze these same vertex labels; it is not an additional test in
\(\mathcal G\).

\begin{lemma}[Completeness of the outer matrix test]
\label{lem:outer-completeness}
If \(\mathcal E\) has an assignment of value at least \(1-\xi\), and
the latent map satisfies
\[
 \Pr_{x\text{ uniform in }\mathcal V,\ a\sim\mu}
       [C(x+a)\ne C(x)]\le p_*,
\]
then
\[
 \val(\mathcal G)\ge 1-k\xi-\frac{p_*}{2}.
\]
\end{lemma}

\begin{proof}
Fix the near-satisfying global bit assignment.  Give every reduced
support a valid answer as follows.  At a singleton use its global bit.
At a retained equation use its first two global bits and set the third
to the unique value satisfying the equation.  These choices depend only
on the support, including its occurrence IDs, and therefore define an
answer for every exact key independently of its origin.

At an orbit representative, evaluate its reduced table on this answer
and apply \(C\).  This defines a vertex labeling.  The equivariance of
\(C\) under \(K\) implies that its restored answer is evaluation of
the actual reduced table followed by \(C\), in agreement with
Equation~\eqref{eq:table-folding}.

With probability at least \(1-k\xi\), every equation in \(U\) is
satisfied by the global assignment.  On such a tuple, the support
answers used for \emph{every} table are projections of the same full
answer \(z_U\in L_U\), whose first coordinate is one.  Consequently
the two sides of Equation~\eqref{eq:matrix-test} are
\[
 C(Pz_U)\quad\text{and}\quad
 C(Pz_U+a(l^\top z_U)).
\]
Since \(z_U\ne0\), the vector \(Pz_U\) is uniform in
\(\mathcal V\), independently of \(a,l\), and \(l^\top z_U\) is
an independent fair bit.  The test therefore fails with probability at
most \(p_*/2\) on these tuples.  Adding the probability of a tuple
containing an unsatisfied equation proves the bound.
\end{proof}
